% Draft abstract and introduction for v7 (2026-10-18). Placeholders in [[ ]] are filled from READOUT46.md.
% Rules: at most four numbers in the abstract; no number in the introduction except "seventeen"; no em-dashes.

\begin{abstract}
A genuine face that was beautified by a camera app, an AR filter or a retouching service still shows the same person doing
the same thing, but a forgery detector may not treat it that way. We score seventeen published detectors on six
beautification sources with one rule and find one failure shared by all of them: a skin-smoothed genuine face is scored as
more forged than the same face untouched, in 89--99\% of pairs for fifteen detectors, and smoothing the skin alone is enough.
We then show where the cue comes from. With architecture, data and recipe fixed, the class that receives smoothed views
during training sets the sign of the error: views on genuine faces only let nearly every smoothed forgery through, views on
forged faces only accuse smoothed genuine faces, and views on both sit in between. From this we derive a training rule,
smoothing-invariant consistency (SIC): one smoothing operator drawn identically for both classes and a penalty on the change
of the forgery logit under it. SIC lowers the false alarm on skin-smoothed genuine faces from [[FA before]]\% to
[[FA after]]\% while leaving untreated detection within the confidence interval of the same model trained without it
([[AUC on Celeb-DF-v2]] on Celeb-DF-v2). We release the six sources, the scorer and the paired reader as a benchmark.
\end{abstract}

\section{Introduction}
Forgery detectors are used for content moderation, remote identity checks and fact-checking. In each setting a genuine face
called fake accuses a real person, and a forgery called real is let through. A large share of the genuine faces reaching
such a detector were beautified before upload: skin smoothing, whitening, eye enlargement and face slimming are default
options in camera apps and retouching services, and colour presets ship with every photo-sharing app. Beautification edits
the pixels a detector inspects but changes neither who is shown nor what the person does.

Prior work measured the effect of beautification on at most six detectors and in one direction at a
time~\cite{libourel2024,saeed2025,concas2025}. We ask three questions. Does the effect hold across the detector families
in use today? Where does it come from? Can it be removed without giving up detection?

\paragraph{Finding.} We score seventeen published detectors, eleven CNN detectors, four CLIP-based detectors and two
generated-image detectors, on six beautification sources with one rule that counts both errors (Fig.~\ref{fig:teaser}).
The failure shared by all of them is skin smoothing: a smoothed genuine face is scored as more forged than the same face
untouched in nearly every pair, smoothing the skin alone with the background untouched suffices, and for the CLIP detectors
the forgery logit grows with the smoothing strength. A radial vignette alone does not produce the shift.

\paragraph{Cause.} We hold architecture, data and recipe fixed and vary only which class receives smoothed views during
training (Fig.~\ref{fig:exposure}). The class that sees smoothing sets the sign of the error. Views on genuine faces only,
the direction taken by fine-tuning on enhanced genuine samples~\cite{saeed2025}, teach that smoothed means genuine, and a
forger who smooths the forgery is let through. Views on forged faces only teach the opposite. Views on both classes, drawn
from the same operator, sit in between. This is shortcut learning~\cite{geirhos2020} under class-dependent
augmentation~\cite{balestriero2022} in a setting where the sign of the resulting error can be measured.

\paragraph{Method.} The cause suggests the fix. Smoothing-invariant consistency (SIC) draws one smoothing operator
identically for every training item and penalises the change of the forgery logit under it (Fig.~\ref{fig:arch}). The
model keeps its three-way output and four operation scores; only the forgery logit is constrained. [[One sentence with the
SIC headline from READOUT46: FA on skin-smoothed genuine faces, miss change, untreated AUC.]]

Our contributions are:
\begin{itemize}
\item A two-sided measurement of seventeen detectors on six beautification sources under one rule, read pair by pair where
originals exist, with a skin-only smoothing test and a dose-response test that isolate the shared cue.
\item A controlled exposure study showing that the class receiving smoothed views sets the sign of the error, including an
evasion route opened by genuine-only views.
\item SIC, a training rule that removes the cue [[at the stated cost]], evaluated on pre-registered bars with two seeds and a
$\lambda = 0$ control, together with the released benchmark.
\end{itemize}
